% ============================================================
%  第 1 章  绪论
% ============================================================

\chapter{绪论}

% ================================================================
\section{研究背景与意义}

\subsection{社交媒体时代的舆情挑战}

随着移动互联网的高速普及与智能手机的全面渗透，微博、B 站、抖音、小红书、知乎等社交媒体平台已成为公众发表观点、传播信息的核心渠道。据中国互联网络信息中心（CNNIC）2025 年统计报告，中国网民规模已突破 11 亿人，其中短视频与社交媒体用户占比超过 90\%~\cite{cnnic2024}。海量用户在这些平台上每天产生数以亿计的评论、弹幕与转发，形成极具参考价值的实时公众情绪镜像。

然而，这一信息富矿同时带来了严峻的舆情管理挑战。其一，数据规模庞大，人工审核难以覆盖；其二，信息传播速度极快，负面舆情往往在数小时内扩散至数百万用户；其三，话题情绪复杂多变，具有强烈的事件驱动特征。针对上述挑战，构建一套能够对多平台评论数据进行实时采集、自动分析与主动预警的系统，具有重要的学术价值与现实意义~\cite{chen2023weibo}。

\subsection{传统舆情系统的局限性}

现有商用舆情监测产品（如鹰眼速读网、清博大数据等）通常采用批处理架构：系统以小时或天为粒度定期从数据平台拉取快照，分析完成后生成静态报告。这一模式存在以下固有缺陷：

\begin{enumerate}
  \item \textbf{时效性不足：}批处理周期通常以小时计，难以在舆情爆发初期即发出预警，错失最佳响应窗口。
  \item \textbf{情感分析粗糙：}多数产品采用基于关键词词典的规则匹配方法，难以理解语境、处理隐性负面情绪及反讽表达。
  \item \textbf{数据孤岛问题：}不同平台的数据采集、分析与展示模块相互独立，全链路数据难以整合追溯。
  \item \textbf{可扩展性受限：}新增监测话题或平台往往需要复杂的系统配置，缺乏灵活性。
\end{enumerate}

\subsection{研究意义}

针对上述痛点，本文设计并实现一套基于 Spark Structured Streaming 的实时舆情分析与热度趋势监测系统。该系统以预训练大语言模型 StructBERT 替代词典规则进行情感判定，以 Delta Lake 湖仓架构替代传统关系型数据库，实现流批一体的数据存储与查询。系统的学术意义在于：（1）探索了预训练中文 NLP 模型在低延迟流处理场景下的工程化部署路径；（2）验证了 Delta Lake Medallion 架构在准实时舆情分析系统中的适用性与可靠性；（3）提出了一种融合用户互动行为与情感极性的多维度热度评分模型。

% ================================================================
\section{国内外研究现状}

\subsection{流式计算框架研究现状}

流式数据处理领域经历了从第一代（Apache Storm~\cite{storm2014}）的纯流式处理，到第二代（Spark Streaming~\cite{zaharia2012discretized}）的微批次处理，再到第三代（Apache Flink~\cite{carbone2015apache}、Spark Structured Streaming~\cite{armbrust2018structured}）的统一流批处理的演进历程。

Zaharia 等人~\cite{zaharia2012discretized}于 2012 年提出的离散化流（DStream）模型将流处理问题转化为一系列小批量 RDD 操作，奠定了 Spark Streaming 的理论基础。Armbrust 等人~\cite{armbrust2018structured}在 2018 年进一步提出 Structured Streaming，将流视为无界关系表，引入 Event-Time 语义与 Watermark 机制，显著提升了处理语义的严谨性与 API 的易用性。相较于 Flink 的原生毫秒级流处理，Structured Streaming 在与 Python 生态（NumPy、PyTorch 等）的集成方面具有更成熟的工程支持，更适合本文所需的"流式调度 + 批量 NLP 推理"混合场景。

\subsection{情感分析技术研究现状}

中文情感分析技术经历了词典规则方法（HowNet 情感词典~\cite{dong2003hownet}）、传统机器学习（SVM、朴素贝叶斯）到深度学习（LSTM~\cite{hochreiter1997lstm}、BERT~\cite{devlin2019bert}）的三个发展阶段。

BERT~\cite{devlin2019bert} 的提出引发了预训练语言模型在情感分析任务上的研究热潮。针对中文特点，阿里达摩院提出的 StructBERT~\cite{wang2020structbert} 在 BERT 基础上引入词序重建与句子顺序预测两类新的预训练任务，显著提升了模型对中文词语结构的理解能力，在 GLUE 基准的多项子任务上超越了 BERT-Large。本系统选用 StructBERT 中文情感分类专项模型（\code{iic/nlp\_structbert\_sentiment-classification\_chinese-base}），实验表明其在社交媒体评论上总体准确率可达 83.5\%（详见第~\ref{sec:api-perf}~节）。

\subsection{数据湖仓技术研究现状}

Delta Lake~\cite{armbrust2020delta} 由 Databricks 团队于 2020 年开源，是基于 Parquet 文件格式构建的开放 ACID 存储层，通过事务日志（Transaction Log）在对象存储上实现 Serializable 级别的隔离性。与传统数据仓库（Hive/Presto）相比，Delta Lake 的核心优势在于：（1）流批统一写入，Structured Streaming 的 \code{writeStream} 可直接写入 Delta 表；（2）Time Travel 支持按版本或时间戳查询历史快照；（3）Schema 演化，允许在不中断下游查询的情况下新增字段。

Medallion 架构（Bronze-Silver-Gold）是 Databricks 推荐的 Delta Lake 数据组织模式~\cite{databricks2021medallion}，已在字节跳动、阿里巴巴等国内大型互联网企业的数据湖建设实践中得到广泛应用。本系统将该架构应用于准实时舆情分析场景，是对其在小规模单机部署条件下适用性的探索性验证。

\subsection{现有研究的不足与本文的切入点}

综合梳理现有工作，可以发现以下不足：其一，多数基于 Spark 的舆情分析研究局限于离线批处理，缺乏对准实时流处理场景的系统性验证；其二，将预训练大模型（BERT 类）嵌入 Structured Streaming \code{foreachBatch} 回调的工程化方案鲜有公开详述；其三，Delta Lake Medallion 架构在中文社交媒体舆情场景中的完整实践案例较少。本文针对上述不足，设计并实现完整的端到端系统，对上述三个方面均作出了有价值的工程实践贡献。

% ================================================================
\section{研究内容与主要工作}

本文的主要研究内容与工作成果概括如下：

\begin{enumerate}
  \item \textbf{多平台社交媒体爬虫工程化实现。}基于 MediaCrawler 框架与 CDP 接管机制，实现了支持微博、B 站、抖音、小红书、知乎、快手、贴吧七平台的关键词评论采集系统，并设计了 JSONL 格式的低耦合 Kafka 数据输出接口。

  \item \textbf{Spark Structured Streaming 实时消费管道。}设计了基于 60~秒微批次触发的 Kafka-Spark 消费管道，实现了 JSON Schema 解析、字段校验与字段异常容错处理。

  \item \textbf{StructBERT 情感分析服务的流处理集成。}提出了懒初始化单例模式+批量推理+异常降级三层设计方案，将 StructBERT Pipeline 嵌入 \code{foreachBatch} 回调，实现单例复用、批量高效推理。

  \item \textbf{多维度热度评分模型设计与实现。}提出了融合点赞数、转发数与情感极性加成的热度评分公式，并设计了基于热度分数与情感标签组合的三级风险评估机制。

  \item \textbf{Delta Lake Bronze-Silver-Gold 三层湖仓架构实现。}设计并实现了原始数据追加（Bronze）、去重精炼 Merge（Silver）与多粒度聚合全量重算（Gold）三层写入方案，保障了数据的 ACID 一致性与多时间粒度查询能力。

  \item \textbf{FastAPI 数据服务与 Vue 3 可视化大屏。}设计了 9 个 RESTful 端点，基于 \code{deltalake} Python 库实现了无 JVM 的 Delta Lake 直读方案；前端基于 Vue~3 Composition API 与 ECharts~6 实现了多话题切换、多粒度趋势图与告警面板的交互式可视化大屏。
\end{enumerate}

% ================================================================
\section{论文组织结构}

本文共分六章，各章内容安排如下：

第一章（绪论）：介绍研究背景、意义及国内外研究现状，阐述本文的研究切入点与主要工作。

第二章（相关技术与理论基础）：系统介绍本文所涉及的核心技术，包括 Apache Kafka、Spark Structured Streaming、Delta Lake、StructBERT、FastAPI 及 Vue 3/ECharts。

第三章（系统需求分析与总体设计）：在需求分析基础上，给出系统五层架构设计、数据流设计、模块划分及关键接口定义。

第四章（系统详细设计与实现）：按模块详细阐述各子系统的设计方案与关键代码实现，重点介绍爬虫模块、流处理管道、情感分析服务、热度计算模块与 Delta Lake 写入逻辑。

第五章（系统测试与实验分析）：描述测试环境与数据集，通过功能测试与性能测试对系统各模块进行验证，并对情感分析准确率、热度模型有效性等核心指标进行综合分析。

第六章（总结与展望）：总结全文工作，客观指出系统的局限性，并对未来改进方向进行展望。
